%topic under study and imply the underlying question. One sentence
%Two sentences: previous research has demonstrated. Provide rationale for new research
%Two sentences: the data, research, and analytical methods used in my study
% major findings and implications and significance of this study
Visual feedback speeds up learners' improvement of pronunciation in a second language. The visual combined with audio allows speakers to see sounds and differences in pronunciation that they are unable to hear. Prior studies have tested different visual methods for improving pronunciation, however, we do not have conclusive understanding of what aspects of the visualizations contributed to improvements. Based on previous work, we created \visowel{}, an interactive vowel chart. Vowel charts provide actionable feedback by directly mapping physical tongue movement onto a chart. We compared \visowel{} with an auditory-only method to explore how learners parse visual and auditory feedback to understand how and why visual feedback is effective for pronunciation improvement. The findings suggest that designers should include explicit anatomical feedback that directly maps onto physical movement for phonetically untrained learners. Furthermore, visual feedback has the potential to motivate more practice since all eight of the participants cited using the visuals as a goal with \visowel{} versus relying on their own judgment with audio alone. Their statements are backed up by all participants practicing words with \visowel{} more than with audio-only. Our results indicate that \visowel{} is effective at providing actionable feedback, demonstrating the potential of visual feedback methods in second language learning.